\begingroup\sloppy
\subsection{Real extended systems, conjugation, stopped transfer}\label{bpN:p11:s-real}

Setting of \S\ref{bpN:p11:s-ext-class}: $I_c$, $S$, $r,s,w$, $\lambda,c$, $\tilde S$. For $0<\varrho\le r$ let
$W^e_\varrho=\{(\theta,I,\tilde x)\in\R^{2n+2k}:\abs{I-I_c}<\varrho,\ \operatorname{dist}_\infty(\tilde x,\tilde S)<\varrho\}=\Dom\varrho s^e\cap\R^{2n+2k}$ (open). For $G\in\Os rs^{e,\R}$ the real Hamiltonian $G^\R$ is $G|_{W^e_r}$ (and $0$ off it). A \emph{$G$-orbit in $W^e_\varrho$ on an interval $J$} is a $y:\R\to\R^{2n+2k}$ with $y(t)\in W^e_\varrho$ and $y'(t)=X^c_{G^\R}(y(t))$ for $t\in J$ (one-sided derivative at endpoints of $J$ if $J\ne\R$; $y$ is defined on $\R$).

\begin{lemma}[Real extended systems]\label{bpN:p11:real-system}
Let $G$ be holomorphic and $1$-periodic in $\theta$ on $\Dom rs^e$, real on $W^e_r$ (finite norm is needed only in (c), (d)), $0<\varrho\le r$.
(a) $G^\R\in C^\infty(W^e_r)$, $G^\R$ is $1$-periodic in $\theta$, and $X^c_{G^\R}(y)=X^c_G(y)\in\R^{2n+2k}$ (complex partials at the real point $y$); $X^c_{G^\R}$ is $C^1$ and $\theta$-periodic on $W^e_r$.
(b) (Energy) If $y$ is a $G$-orbit in $W^e_\varrho$ on $J$ then $G^\R\circ y$ is constant on $J$.
(c) For $u\in\Os\varrho s^e$ with finite norm and $y\in W^e_\varrho$: $\abs{u(y)}\le\wn u\varrho s$ and $\sum_j\abs{\partial_{\theta_j}u(y)}\le\wn u\varrho s/(es)$.
(d) If $P_\Lambda g=g$ ($g\in\Os\varrho s^e$ finite norm, $\Lambda\subset\Z^n$) then $\Pi_\Lambda^\perp\partial_\theta g=0$ on $W^e_\varrho$.
\end{lemma}
\begin{proof}
\proofstep{(a)}\label{stepN:p11:rs-a} Lemma~\ref{bpN:p7:real-system}(a),(b),(c) with $\Dom rs^e\cap\R^{2n+2k}=W^e_r$ and \Kref{bp:p1:hamiltonian-real}: the proof only uses holomorphy on an open set and reality at the real points, not the structure of the variables.
\proofstep{(b)}\label{stepN:p11:rs-b} $G^\R$ is differentiable and $\frac d{dt}G^\R(y(t))=\nabla G^\R(y)\cdot\Omega_c\nabla G^\R(y)=0$ ($\Omega_c^\top=-\Omega_c$); a function with zero derivative on an interval is constant (mean value inequality, \Kref{bp:p1:ext-mvt}).
\proofstep{(c),(d)}\label{stepN:p11:rs-cd} Lemma~\ref{bpN:p7:theta-derivative} and Lemma~\ref{bpN:p7:resonant-gradient}: both are statements about Fourier coefficients in $\theta$ at fixed $\hat I$ and use Lemma~\ref{bpN:p11:ext-part4} only.
\end{proof}

\begin{lemma}[Conjugation]\label{bpN:p11:conj}
Let $H\in\Os rs^{e,\R}$, $\tilde H\in\Os{\tilde r}{\tilde s}^{e,\R}$ ($\tilde r\le r$, $\tilde s\le s$ in the same extended structure), $\Psi:\Dom{\tilde r}{\tilde s}^e\to\Dom rs^e$ holomorphic, real at real points, with $D\Psi\,\Omega_c\,D\Psi^\top=\Omega_c$ and $\tilde H=H\circ\Psi$. If $\tilde y$ is an $\tilde H$-orbit in $W^e_{\tilde r}$ on $J$ then $\Psi\circ\tilde y$ is an $H$-orbit in $W^e_r$ on $J$.
\end{lemma}
\begin{proof}
\proofstep{Chain rule}\label{stepN:p11:cj} As Lemma~\ref{bpN:p7:conjugation}: $Y=D\Psi(\tilde y)$ is real, $\nabla\tilde H^\R=Y^\top\nabla H^\R\circ\Psi$, so $\frac d{dt}\Psi(\tilde y)=Y\Omega_cY^\top\nabla H^\R(\Psi\tilde y)=X^c_{H^\R}(\Psi\tilde y)$; $\Psi$ maps $W^e_{\tilde r}$ into $W^e_r$.
\end{proof}

\begin{lemma}[Real preimages of a near-identity map]\label{bpN:p11:preimage}
Let $\Psi:\Dom{\tilde r}{\tilde s}^e\to\C^{2n+2k}$ be holomorphic, real at real points, $G=\Psi-\mathrm{id}$, $\abs{G_\theta}\le d_\theta$, $\abs{G_I}\le d_I$, $\abs{G_{\tilde x}}\le d_x$ on $\Dom{\tilde r}{\tilde s}^e$. Let $0<\varrho'\le\tilde r$, $0<\sigma'\le2\tilde s$ and $\kappa_*=\max(d_\theta/\sigma',d_I/\varrho',d_x/\varrho')\le\frac14$. Then every $y=(\theta_y,I_c,\tilde x_y)$ with $\theta_y\in\R^n$, $\tilde x_y\in\tilde S$ equals $\Psi(x)$ for some $x\in W^e_{\tilde r}$ with $\abs{\theta_x-\theta_y}\le d_\theta$, $\abs{I_x-I_c}\le d_I$, $\abs{\tilde x_x-\tilde x_y}\le d_x$.
\end{lemma}
\begin{proof}
Proof of Lemma~\ref{bpN:p7:preimage} with $\nu(a,b,\xi)=\max(\abs a/\sigma',\abs b/\varrho',\abs\xi/\varrho')$: $X=\{x\in\R^{2n+2k}:\nu(x-y)\le\kappa_*\}\subset W^e_{\tilde r}$ (since $\tilde x_y\in\tilde S$ and $\kappa_*\varrho'\le\varrho'/4<\tilde r$); $T(x)=y-G(x)$ maps $X$ to itself; the complex $\nu$-ball of radius $\frac12$ about $x\in X$ lies in $\Dom{\tilde r}{\tilde s}^e$ (imaginary parts $<\sigma'/2\le\tilde s$, $\operatorname{dist}(\tilde x',\tilde S)<\frac{\varrho'}4+\frac{\varrho'}2<\tilde r$), so Lemma~\ref{bpN:p4:general-cauchy} gives $\abs{DG}_\nu\le2\kappa_*\le\frac12$ and $T$ is a contraction on the convex set $X$ (mean value inequality); Banach's fixed point theorem (\ref{bpN:p7:ext-banach}) gives $x$.
\end{proof}

\begin{lemma}[Stopped transfer through a symplectic map]\label{bpN:p11:transfer}
Let $H,\tilde H,\Psi$ be as in Lemma~\ref{bpN:p11:conj}, $\abs\cdot_*\ge\abs\cdot$ a norm on $\R^n$, and assume
\begin{enumerate}[label=(T\arabic*)]
\item $\abs{\Pi_I\Psi(x)-I_x}_*\le d_*$ and $\abs{\Pi_{\tilde x}\Psi(x)-\tilde x_x}_\infty\le d_x$ for all $x\in W^e_{\tilde r}$, with $d_x<\tilde r$;
\item every $y=(\theta_y,I_c,\tilde x_y)$ ($\theta_y\in\R^n$, $\tilde x_y\in\tilde S$) is $\Psi(x)$ for some $x\in W^e_{\tilde r}$;
\item there are $\varrho>0$ with $d_*+\varrho<\tilde r$ and $T\in(0,\infty]$ such that every $\tilde H$-orbit $\tilde y$ in $W^e_{\tilde r}$ on an interval $\tilde J\ni0$ with $\abs{\tilde I(0)-I_c}_*\le d_*$ satisfies $\abs{\tilde I(t)-\tilde I(0)}_*\le\varrho$ for $t\in\tilde J$, $\abs t\le T$.
\end{enumerate}
Let $F:\R^{2n+2k}\to\R^{2n+2k}$ be any function with $F=X^c_{H^\R}$ on $W^e_r$, $J\ni0$ an interval, $y:\R\to\R^{2n+2k}$ differentiable on $J$ with $y'=F(y)$ on $J$, $y(0)=(\theta_0,I_c,\tilde x_0)$, $\theta_0\in\R^n$, $\tilde x_0\in\tilde S$. Let $t_1\in J$, $\abs{t_1}\le T$, and suppose $\tilde x(u)\in\tilde S$ for all $u$ between $0$ and $t_1$ ($\tilde x$ the slow components of $y$). Then $y(u)\in W^e_r$ and $\abs{I(u)-I_c}_*\le\varrho+2d_*$ for all $u$ between $0$ and $t_1$.
\source{Ours; Lemma~\ref{bpN:p7:transfer} and Corollary~\ref{bpN:p7:stopped-extension} combined, with the new continuation argument for the slow variables (the normal-form orbit may leave $W^e_{\tilde r}$ through the slow variables; the stopped hypothesis prevents it).}
\end{lemma}
\begin{proof}
Take $t_1>0$; for $t_1<0$ apply the result to $u\mapsto y(-u)$ and $-H$, $-\tilde H$ (Lemma~\ref{bpN:p3:gen-symmetry}(b): $X^c_{-G}=-X^c_G$), and for $t_1=0$ there is nothing to prove.
\proofstep{A normal-form orbit}\label{stepN:p11:tr-orbit} By (T2) choose $x\in W^e_{\tilde r}$ with $\Psi(x)=y(0)$; by (T1) $\abs{I_x-I_c}_*\le d_*$. Let $\tilde y$ be the maximal $\tilde H$-orbit in $W^e_{\tilde r}$ through $x$ on $(T_-,T_+)$ (Lemma~\ref{bpN:p11:real-system}(a), Theorem~\ref{bpN:p3:gen-maximal}), $T_-<0<T_+$.
\proofstep{The good set}\label{stepN:p11:tr-E} Let $E$ be the set of $\tau\in[0,t_1]$ such that $[0,\tau]\subset(T_-,T_+)$, $y(u)\in W^e_r$ for $u\in[0,\tau]$ and $y=\Psi\circ\tilde y$ on $[0,\tau]$. $0\in E$.
\proofstep{$E$ is open in $[0,t_1]$}\label{stepN:p11:tr-open} Let $\tau\in E$, $\tau<t_1$. $\Psi\circ\tilde y$ is an $H$-orbit in $W^e_r$ on $(T_-,T_+)$ (Lemma~\ref{bpN:p11:conj}). \emph{Claim $\tau<T_+$:} if $\tau\ge T_+$ then (as $[0,\tau]\subset(T_-,T_+)$) this is impossible; so $\tau<T_+$. $y(\tau)=\Psi(\tilde y(\tau))\in W^e_r$, open, and $y$ is continuous, so $y(u)\in W^e_r$ for $u$ near $\tau$, where $F=X^c_{H^\R}$; $y$ and $\Psi\circ\tilde y$ are two solutions of $z'=X^c_{H^\R}(z)$ on a neighbourhood of $\tau$ in $W^e_r$ with the same value at $\tau$, hence equal there (Theorem~\ref{bpN:p3:gen-uniqueness}; $X^c_{H^\R}$ is $C^1$). So $[\tau,\tau+\epsilon)\subset E$.
\proofstep{A priori bounds on $E$}\label{stepN:p11:tr-bounds} Let $\tau\in E$. Hypothesis (T3) applies to $\tilde y$ on $\tilde J=[0,\tau]$ extended to $(T_-,T_+)\cap\{\ldots\}$: for $u\in[0,\tau]$, $\abs{\tilde I(u)-\tilde I(0)}_*\le\varrho$; $\tilde I(0)=I_x$ with $\abs{I_x-I_c}_*\le d_*$, so $\tilde I(u)\in\{\abs{I-I_c}<\tilde r\}$ (as $\varrho+d_*<\tilde r$ and $\abs\cdot\le\abs\cdot_*$). By (T1), $\abs{\tilde x_{\tilde y}(u)-\tilde x(u)}\le d_x$ where $\tilde x(u)\in\tilde S$ is the slow component of $y(u)=\Psi(\tilde y(u))$ (assumption). Hence $\operatorname{dist}(\tilde x_{\tilde y}(u),\tilde S)\le d_x<\tilde r$.
\proofstep{$E$ is closed}\label{stepN:p11:tr-closed} Let $\tau^*=\sup E$. The image $\tilde y([0,\tau^*))$ lies in $K=\{\theta\in\R^n\}\times\{\abs{I-I_x}_*\le\varrho\}\times\{\tilde x:\operatorname{dist}(\tilde x,\tilde x([0,t_1]))\le d_x\}$, a set that is compact modulo the $\Z^n$-translations of $\theta$ ($\tilde x([0,t_1])$ is the continuous image of $[0,t_1]$, hence compact), and which is contained in $W^e_{\tilde r}$ by Step~\ref{stepN:p11:tr-bounds}; the field $X^c_{\tilde H^\R}$ is periodic, so Corollary~\ref{bpN:p3:gen-continuation-periodic} gives $T_+>\tau^*$. By continuity $\tilde y(\tau^*)$, $y(\tau^*)=\Psi(\tilde y(\tau^*))\in W^e_r$ and $y=\Psi\circ\tilde y$ on $[0,\tau^*]$, i.e.\ $\tau^*\in E$.
\proofstep{Conclusion}\label{stepN:p11:tr-concl} $E$ is nonempty, open and closed in $[0,t_1]$, so $E=[0,t_1]$. For $u\in[0,t_1]$: $I(u)-I_c=(\Pi_I\Psi(\tilde y(u))-\tilde I(u))+(\tilde I(u)-\tilde I(0))+(I_x-I_c)$ (using $\tilde I(0)=I_x$), of size $\le d_*+\varrho+d_*$ in $\abs\cdot_*$ by (T1), (T3), and $\abs{I_x-I_c}_*\le d_*$.
\end{proof}

\subsection{Nonresonant stability (Proposition 1$^e$)}\label{bpN:p11:s-nr}

\begin{proposition}[Nonresonant stability in the extended phase space]\label{bpN:p11:nonresonant}
Let $r,s,K_a,\alpha>0$, $L\ge0$, $p>1$, $q=p/(p-1)$, $\kappa=2\pi K_as\ge6$; $h$ holomorphic on $V_r(I_c)$, real on real points; $f\in\Os rs^{e,\R}$, $\varepsilon=\wn frs$; $H=h+f$. Assume \ref{bpN:p6:cond-nr} for $\Lambda=\{0\}$ ($\abs{k\cdot\omega(I_c)}\ge\alpha$ for $0<\abs k_1\le K_a$), \ref{bpN:p6:cond-L} with constant $L$ on $V_r(I_c)$, $K_aLr\le\alpha/p$, (SL), and $\varepsilon\le\alpha r/(2^7qK_a)$. Put $T_{\rm nr}=\frac{sr}{5\varepsilon}e^{\kappa/6}$. Then for every $F$ with $F=X^c_{H^\R}$ on $W^e_r$, every $J\ni0$ and $y$ with $y'=F(y)$ on $J$, $y(0)=(\theta_0,I_c,\tilde x_0)$, $\tilde x_0\in\tilde S$, and every $t_1\in J$ with $\abs{t_1}\le T_{\rm nr}$ and $\tilde x(u)\in\tilde S$ between $0$ and $t_1$: $\abs{I(u)-I_c}\le r/2$ for all $u$ between $0$ and $t_1$.
\source{Pos93 p.~196 (Proposition 1), Proposition~\ref{bpN:p7:nonresonant}.}
\end{proposition}
\begin{proof}
\proofstep{Normal form}\label{stepN:p11:nr-nf} Theorem~\ref{bpN:p11:nf} with $\Lambda=\{0\}$ ($g_0=\hat f_0$ is $\theta$-independent): $\Psi:\Dom{r/2}{s/6}^e\to\Dom rs^e$, $H\circ\Psi=h+g+f_*$, $\eta=\wn{f_*}{r/2}{s/6}\le e^{-\kappa/6}\varepsilon$, $\abs{\Pi_I\Psi-I}\le d_I:=qK_a\varepsilon/\alpha\le r/2^7$, $\abs{\Pi_{\tilde x}\Psi-\tilde x}\le r/2^{10}$, $\abs{\Pi_\theta\Psi-\theta}\le d_\theta=10q\varepsilon/(2\pi\alpha r)$.
\proofstep{(T1), (T2)}\label{stepN:p11:nr-t12} With $\tilde r=r/2$, $\tilde s=s/6$, $\abs\cdot_*=\abs\cdot$, $d_*=d_I$, $d_x=r/2^{10}$: (T1) is the displacement. (T2): Lemma~\ref{bpN:p11:preimage} with $\varrho'=r/2$, $\sigma'=s/3$: $d_I/\varrho'\le\frac1{64}$, $d_x/\varrho'\le2^{-9}$, $d_\theta/\sigma'=\frac{30q\varepsilon}{2\pi\alpha rs}\le\frac{30}{2\pi2^7K_as}=\frac{30}{2^7\kappa}\le\frac{30}{768}<\frac14$.
\proofstep{(T3)}\label{stepN:p11:nr-t3} Take $\varrho=\frac{15}{32}r$: $d_*+\varrho\le\frac{61}{128}r<\tilde r$. Let $\tilde y$ be an $\tilde H$-orbit in $W^e_{\tilde r}$ on $\tilde J$, $\abs{\tilde I(0)-I_c}\le d_*$, and suppose $\abs{\tilde I(t_1')-\tilde I(0)}>\varrho$ with $\abs{t_1'}\le T_{\rm nr}$. Lemma~\ref{bpN:p7:exit-ball} gives $t^*$, $\abs{t^*}\le\abs{t_1'}$, $\abs{\tilde I(t^*)-\tilde I(0)}=\varrho$ and $\tilde I(\tau)$ in the closed $\varrho$-ball about $\tilde I(0)$ before $t^*$ (inside $\abs{I-I_c}<\tilde r$). On this stretch $\dot{\tilde I}=-\partial_\theta\tilde H^\R=-\partial_\theta f_*(\tilde y)$ ($h$, $g$ are $\theta$-independent; the slow components do not enter the $I$-equation), so by Lemma~\ref{bpN:p11:real-system}(c) $\abs{\dot{\tilde I}}\le\eta/(e\tilde s)=6\eta/(es)$ and $\varrho\le\frac{6\eta\abs{t^*}}{es}\le\frac{6e^{-\kappa/6}\varepsilon}{es}\cdot\frac{sre^{\kappa/6}}{5\varepsilon}=\frac6{5e}r<0.442r<\frac{15}{32}r$: contradiction. (If $\varepsilon=0$, $\eta=0$.)
\proofstep{Transfer}\label{stepN:p11:nr-concl} Lemma~\ref{bpN:p11:transfer} gives $\abs{I(u)-I_c}\le\frac{15}{32}r+\frac2{2^7}r=\frac{31}{64}r\le\frac r2$.
\end{proof}

\subsection{The resonant case: an $\omega$-free energy error}\label{bpN:p11:s-energy}

\begin{lemma}[Energy confinement near a resonance, convex case, no bound on $\omega$]\label{bpN:p11:energy-free}
Let $\tilde r,\tilde s>0$, $I_c\in\R^n$, $\Lambda\subset\Z^n$ with Euclidean projection $\Pi=\Pi_\Lambda$, $\Pi^\perp=1-\Pi$. Let $h$ be holomorphic on $V_{\tilde r}(I_c)$, real on real points, with, on $U=\{I\in\R^n:\abs{I-I_c}<\tilde r\}$, $m\norm\xi_2^2\le Q(I)[\xi,\xi]$ and $\norm{Q(I)}_{2\to2}\le M$ ($Q=D^2h_\R$, $0<m\le M$). Let $g,f_*\in\Os{\tilde r}{\tilde s}^{e,\R}$ with finite norms, $P_\Lambda g=g$, $\zeta_g=\wn g{\tilde r}{\tilde s}$, $\eta=\wn{f_*}{\tilde r}{\tilde s}$, and let $\varepsilon>0$ with
\[\zeta_g\le\tfrac54\varepsilon,\qquad\eta\le\varepsilon/e,\qquad\varepsilon\le m\tilde r^2/64,\]
\[16kc\tilde s\le\tilde r\quad(\text{this is }\textup{(SL)}_{\tilde r/2,\tilde s/2}\text{ for the extended structure}).\]
Let $\delta\in(0,m\tilde r/4]$ with $\norm{\Pi\omega(I_c)}_2\le\delta$, $\rho'=\delta/(4M)$ ($\le\tilde r/16$), $\tilde H=h+g+f_*$, $T_*=\dfrac{\tilde s\,m\,\tilde r}{8M\eta}$ ($=+\infty$ if $\eta=0$). Then for every $\tilde H$-orbit $\tilde y$ in $W^e_{\tilde r}$ on an interval $\tilde J\ni0$ with $\norm{\tilde I(0)-I_c}_2<\rho'$:
\[\norm{\tilde I(t)-\tilde I(0)}_2\le\tilde r-\rho'\qquad\text{for }t\in\tilde J,\ \abs t\le T_* .\]
\source{Pos93 pp.~197--198 (Lemma 1 and its proof: $\frac m3\norm{\Delta I}^2\le2\zeta+\frac5{2m}\delta^2+T\norm{\omega_0}\frac{\abs{f_*}}{es}$), Lemma~\ref{bpN:p7:no-exit}. \emph{New:} the term $\norm{\omega_0}\norm{\Pi^\perp\Delta I}$ is estimated by the exact identity \eqref{bpN:p11:eq:lie} below, whence no bound on $\norm\omega$ and the time $T_*$ is independent of $\omega_0$ (Remark~\ref{bpN:p11:omega-unbounded}).}
\end{lemma}
\begin{proof}
Suppose some $t_1\in\tilde J$, $\abs{t_1}\le T_*$, violates the claim. Lemma~\ref{bpN:p7:exit-ball} ($\abs\cdot_*=\norm\cdot_2$) gives $t^*$, $\abs{t^*}\le T_*$, with $\norm{\tilde I(\tau)-\tilde I(0)}_2\le\Delta:=\tilde r-\rho'$ for $\tau$ between $0$ and $t^*$ and equality at $t^*$. Put $\omega_0=\omega(\tilde I(0))$, $\Delta I=\tilde I(t^*)-\tilde I(0)$, $I_{s'}=\tilde I(0)+s'\Delta I$.
\proofstep{Admissibility}\label{stepN:p11:ef-adm} The ball $\{\norm{I-\tilde I(0)}_2\le\Delta\}\subset U$ (it is convex and $\abs{I-I_c}\le\norm{I-\tilde I(0)}_2+\norm{\tilde I(0)-I_c}_2<\Delta+\rho'=\tilde r$). The mean value inequality for $\omega$ ($D\omega=Q$, $\norm Q\le M$): $\norm{\omega_0-\omega(I_c)}_2\le M\rho'=\delta/4$, so $\norm{\Pi\omega_0}_2\le\frac54\delta$. Also $\rho'\le\frac{m\tilde r}{16M}\le\frac{\tilde r}{16}$, so $\Delta\ge\frac{15}{16}\tilde r$.
\proofstep{Energy inequality}\label{stepN:p11:ef-en} As in Lemma~\ref{bpN:p7:energy-core}(a),(c),(d) (the energy is Lemma~\ref{bpN:p11:real-system}(b); the Taylor formula with Lagrange remainder along $I_{s'}\subset U$, convexity $Q\ge m$): $\frac m2\Delta^2\le A+\abs{\ip{\Pi\omega_0}{\Delta I}}+\abs{\ip{\omega_0}{\Pi^\perp\Delta I}}$, with $A=\abs{(g+f_*)(\tilde y(t^*))-(g+f_*)(\tilde y(0))}\le2(\zeta_g+\eta)$. Here $\abs{\ip{\Pi\omega_0}{\Delta I}}\le\frac54\delta\Delta\le\frac5{16}m\tilde r\Delta\le\frac5{16}\cdot\frac{16}{15}m\Delta^2=\frac13m\Delta^2$. Hence
\[\frac m6\Delta^2\le A+C,\qquad C:=\abs{\ip{\omega_0}{\Pi^\perp\Delta I}} .\]
\proofstep{The identity}\label{stepN:p11:ef-id} Along $\tilde y$, with $\tilde H=h+g+f_*$ and $h$ independent of $(\theta,\tilde x)$, the function $\mathsf F(\tau)=f_*(\tilde y(\tau))$ has derivative $\mathsf F'=Df_*\cdot X^c_{\tilde H^\R}=\{f_*,\tilde H\}_c=\{f_*,h\}_c+\{f_*,g\}_c=\partial_\theta f_*\cdot\omega(\tilde I)+\{f_*,g\}_c$ (the bracket $\{f_*,f_*\}_c=0$; Lemma~\ref{bpN:p11:real-system}(a) for the field). That is
\begin{equation}\label{bpN:p11:eq:lie}\partial_\theta f_*\cdot\omega(\tilde I(\tau))=\mathsf F'(\tau)-\{f_*,g\}_c(\tilde y(\tau)).\end{equation}
\proofstep{Estimate of $C$}\label{stepN:p11:ef-C} $\dot{\tilde I}=-\partial_\theta(g+f_*)=-\partial_\theta f_*-\partial_\theta g$ and $\Pi^\perp\partial_\theta g=0$ (Lemma~\ref{bpN:p11:real-system}(d)), so $\Pi^\perp\Delta I=-\int_0^{t^*}\Pi^\perp\partial_\theta f_*(\tilde y)\,d\tau$ and, $\Pi$ being self-adjoint, $\ip{\omega_0}{\Pi^\perp\Delta I}=-\int\ip{\omega_0}{\partial_\theta f_*}+\int\ip{\Pi\omega_0}{\partial_\theta f_*}$. Insert $\ip{\omega_0}{\partial_\theta f_*}=\ip{\omega(\tilde I)}{\partial_\theta f_*}+\ip{\omega_0-\omega(\tilde I)}{\partial_\theta f_*}$ and \eqref{bpN:p11:eq:lie}:
\[C\le\abs{\mathsf F(t^*)-\mathsf F(0)}+\abs{t^*}\Bigl(\sup\abs{\{f_*,g\}_c}+M\Delta\sup\norm{\partial_\theta f_*}_2+\tfrac54\delta\sup\norm{\partial_\theta f_*}_2\Bigr).\]
Now $\abs{\mathsf F(t^*)-\mathsf F(0)}\le2\eta$. By Lemma~\ref{bpN:p11:real-system}(c) $\norm{\partial_\theta f_*}_2\le\eta/(e\tilde s)$. By Lemma~\ref{bpN:p11:bracket} with $(u,v)=(f_*,g)$, losses $(\rho,\sigma)=(\tilde r/2,\tilde s/2)$ ($\textup{(SL)}_{\rho,\sigma}$ is $16kc\tilde s/2\le\tilde r/2$, the hypothesis) and $r_0=r=\tilde r$: $\sup\abs{\{f_*,g\}_c}\le\frac1{\rho\sigma}\eta\zeta_g=\frac4{\tilde r\tilde s}\eta\zeta_g\le\frac{5\eta\varepsilon}{\tilde r\tilde s}\le\frac{5m\tilde r}{64\tilde s}\eta$. With $M\Delta\le M\tilde r$, $\frac54\delta\le\frac5{16}m\tilde r$ and $\frac1e<0.368$:
\[C\le2\eta+\frac{\abs{t^*}\eta}{\tilde s}\bigl(0.0782\,m\tilde r+0.368\,M\tilde r+0.115\,m\tilde r\bigr)\le2\eta+0.561\,\frac{M\tilde r\eta\abs{t^*}}{\tilde s}\le2\eta+0.0702\,m\tilde r^2,\]
using $\abs{t^*}\le T_*=\tilde sm\tilde r/(8M\eta)$ and $m\le M$.
\proofstep{Contradiction}\label{stepN:p11:ef-contra} $A\le2(\frac54+\frac1e)\varepsilon\le3.24\varepsilon\le0.0507\,m\tilde r^2$ and $2\eta\le\frac{2\varepsilon}e\le0.0115\,m\tilde r^2$ ($\varepsilon\le m\tilde r^2/64$); so $A+C\le0.1324\,m\tilde r^2$. But $\frac m6\Delta^2\ge\frac m6\cdot\frac{225}{256}\tilde r^2>0.1464\,m\tilde r^2$ (\ref{bpN:p11:ext-num}). Contradiction.
\end{proof}

\begin{proposition}[Resonant stability in the extended phase space, convex case]\label{bpN:p11:resonant}
Let $r,s,K_a,\alpha,\delta>0$, $1<p\le\frac32$, $q=p/(p-1)$, $\kappa=2\pi K_as\ge6$, $\Lambda\subset\Z^n$ a subgroup, $0<m\le M$. Let $h$ be holomorphic on $V_r(I_c)$, real on real points, with $\norm{\partial_I^2h}_{2\to2}\le M$ on $V_r(I_c)$ and $m$-convex on $U_r(I_c)$; $f\in\Os rs^{e,\R}$, $\varepsilon=\wn frs$; $H=h+f$. Assume \ref{bpN:p6:cond-nr} modulo $\Lambda$ at $I_c$ ($\abs{k\cdot\omega(I_c)}\ge\alpha$ for $k\in Z_{K_a}\setminus\Lambda$), \ref{bpN:p6:cond-L} with $L=\sqrt nM$ on $V_r(I_c)$ (implied by the Hessian bound, Lemma~\ref{bpN:p6:lipschitz}), $\norm{\Pi_\Lambda\omega(I_c)}_2\le\delta\le mr/8$, (SL), $\varepsilon\le\frac pq\cdot\frac{mr^2}{2^7}$ and $r\le\alpha/(p\sqrt n\,MK_a)$. Put $T_\Lambda=\dfrac{s\,m\,r}{100\,M\,\varepsilon}e^{\kappa/6}$. Then for every $F$ with $F=X^c_{H^\R}$ on $W^e_r$, $J\ni0$, $y$ with $y'=F(y)$ on $J$, $y(0)=(\theta_0,I_c,\tilde x_0)$, $\tilde x_0\in\tilde S$, and $t_1\in J$ with $\abs{t_1}\le T_\Lambda$, $\tilde x(u)\in\tilde S$ between $0$ and $t_1$: $\norm{I(u)-I_c}_2\le r/2$ (so $\abs{I(u)-I_c}\le r/2$) for $u$ between $0$ and $t_1$.
\source{Pos93 pp.~199--200 (Proposition 2) in the convex case, with Lemma 1 (p.~197) in place of Lemma 2, so that the hypotheses $r\le4\ell/m$, $Mr\le\ell/2$ (Remark~\ref{bpN:p7:errata}(1)) are not needed; the time bound is new (Remark~\ref{bpN:p11:omega-unbounded}).}
\end{proposition}
\begin{proof}
\proofstep{Reductions}\label{stepN:p11:rs-red} $\frac pq\le\frac12$, so $\varepsilon\le mr^2/2^8$. By $r\le\alpha/(p\sqrt nMK_a)$ and $m\le M$: $\varepsilon\le\frac{p\,Mr^2}{q2^7}\le\frac{\alpha r}{2^7qK_a}$, i.e.\ hypothesis (i) of Theorem~\ref{bpN:p11:nf}. $\Lambda$ is a subgroup, so $\Lambda=-\Lambda$.
\proofstep{Normal form}\label{stepN:p11:rs-nf} Theorem~\ref{bpN:p11:nf}: $\Psi$, $g$, $f_*$ on $\Dom{r/2}{s/6}^e$, $\wn g{r/2}{s/6}\le\wn{g_0}rs+16c_1\varepsilon^2\le\frac54\varepsilon$ (as $16c_1\varepsilon=\frac{32qK_a\varepsilon}{\alpha r}\le\frac14$), $\eta=\wn{f_*}{}{}\le e^{-\kappa/6}\varepsilon\le\varepsilon/e$ ($\kappa\ge6$). Put $\tilde r=r/2$, $\tilde s=s/6$, $\delta'=mr/8=m\tilde r/4\ge\delta$, $\rho'=\delta'/(4M)=\frac{mr}{32M}$. Then $\varepsilon\le mr^2/2^8=m\tilde r^2/64$, and the slow hypothesis of Lemma~\ref{bpN:p11:energy-free}, $16kc\tilde s\le\tilde r$, follows from $c\le r/(32ks)$: $16kc\frac s6\le\frac{r}{12}<\frac r2$.
\proofstep{(T1), (T2) in $\abs\cdot_*=\norm\cdot_2$}\label{stepN:p11:rs-t12} $\norm{\Pi_I\Psi-I}_2\le\sqrt n\,qK_a\varepsilon/\alpha\le\frac{q\varepsilon}{pMr}\le\frac{mr}{2^7M}=\rho'/4=:d_*$ (the bound $\alpha\ge p\sqrt nMK_ar$, then $\varepsilon\le\frac pq\frac{mr^2}{2^7}$), $d_x=r/2^{10}$; (T2) is Lemma~\ref{bpN:p11:preimage} as in Step~\ref{stepN:p11:nr-t12} ($\kappa\ge6$).
\proofstep{(T3)}\label{stepN:p11:rs-t3} $\varrho=\tilde r-\rho'$, $T=\tilde T_*$ of Lemma~\ref{bpN:p11:energy-free}. An $\tilde H$-orbit with $\norm{\tilde I(0)-I_c}_2\le d_*<\rho'$ is admissible, so $\norm{\tilde I(t)-\tilde I(0)}_2\le\varrho$ for $\abs t\le\tilde T_*=\frac{\tilde sm\tilde r}{8M\eta}=\frac{s\,m\,r}{96M\eta}\ge\frac{smr\,e^{\kappa/6}}{96M\varepsilon}\ge T_\Lambda$ ($\eta\le e^{-\kappa/6}\varepsilon$). Also $d_*+\varrho<\tilde r$ and $\abs\cdot\le\norm\cdot_2$.
\proofstep{Conclusion}\label{stepN:p11:rs-concl} Lemma~\ref{bpN:p11:transfer} (with $T=T_\Lambda$): $\norm{I(u)-I_c}_2\le\varrho+2d_*=\tilde r-\rho'+\rho'/2<r/2$.
\end{proof}

\subsection{The stopped block theorem}\label{bpN:p11:s-block}

\begin{lemma}[Hessian and sizes]\label{bpN:p11:hessian}
Let $h$ be as in Theorem~\ref{bpN:p0:partial}, $\sigma\in\{1,-1\}$ as in \texttt{HasDefiniteFastPart}. For $z\in G'_\rho$: the symmetric matrix $Q_{\C}(z)=(D^2h(z)[e_i,e_j])_{ij}$ satisfies $\norm{Q_{\C}(z)}_{2\to2}\le M_{\C}$. For real $I\in G'_\rho$: $Q(I)=D^2h_\R(I)=Q_{\C}(I)$ (Lemma~\ref{bpN:p1:holo-real}), $m\norm\xi_2^2\le\sigma Q(I)[\xi,\xi]$ and $\abs{Q(I)[\xi,\xi]}\le M_{\C}\norm\xi_2^2$; in particular $m\le M_{\C}$, and $\abs{\omega(I)}_\infty\ge g$. For $h_1=h/2\pi$ the same holds with $(m_1,M_1,g_1)$, and $\sigma h_1$ is $m_1$-convex.
\end{lemma}
\begin{proof}
\proofstep{Operator norm}\label{stepN:p11:hs-op} $\abs{D^2h(z)[u,v]}\le M_{\C}\abs u_\infty\abs v_\infty\le M_{\C}\norm u_2\norm v_2$, and $\norm{Q_{\C}u}_2=\sup_{\norm v_2=1}\abs{\sum_iv_i(Q_{\C}u)_i}=\sup\abs{D^2h(z)[u,v]}$ (take $v\to\bar v$). The real Hessian is the restriction (Lemma~\ref{bpN:p1:holo-real}); the bounds follow from \texttt{HasDefiniteFastPart} (with $\sigma=1$) and $\abs\xi_\infty\le\norm\xi_2$.
\end{proof}

\begin{definition}[Sets and the extended system for the block theorem]\label{bpN:p11:block-setting}
Data and constants of Definition~\ref{bpN:p11:constants}, $\mu>0$ with $\varepsilon=\varepsilon(\mu)\le\varepsilon_0$ and $d=\mu^{b_n}\le\rho$. Put $S=\operatorname{core}(D,3d/4)\subset\mathrm{Ph}_k=\R^{2k}$, $w=d/4$. The \emph{perturbation pair} is $(\breve h,\breve f)=(\sigma h_1,\sigma f_1)$, where $f_1(\theta,I,x)=\frac\mu{2\pi}P\bigl((2\pi\theta,q),(I,p)\bigr)$ for $x=(q,p)$. For $I_c\in G$ and block data $(r,\lambda=w/r,c=2\pi/\lambda^2)$ put $W^e_r(I_c)$ as in \S\ref{bpN:p11:s-real}. The physical real set is $\mathcal W(I_c,r)=\{(\theta,I;x):\abs{I-I_c}<r,\ \operatorname{dist}_\infty(x,S)<w\}$.
\end{definition}

\begin{lemma}[Analyticity, norm and field on the block domain]\label{bpN:p11:domain}
In Definition~\ref{bpN:p11:block-setting}, let $I_c\in G$, $0<r\le r_0$. (a) $\breve h$ is holomorphic on $V_r(I_c)\subset G'_\rho$, real on real points, with $\norm{\partial_I^2\breve h}_{2\to2}\le M_1$ there, $m_1$-convex on $U_r(I_c)$ with $\abs{\omega(I)}_\infty\ge g_1$. (b) $\breve f\in\Os rs^{e,\R}$ for the extended structure with base $(I_c,S)$, widths $(r,s,w)$, and $\wn{\breve f}rs\le\varepsilon(\mu)=\kappa_C\mu/(2\pi)$. (c) On $\mathcal W(I_c,r)$ (open) the frozen Hamiltonian $\hat H_\mu=\frac1{2\pi}H_\mu\circ\Theta^{-1}$ equals $\operatorname{Re}(h_1(I)+f_1)$ and is real analytic; hence $\pm X^{2\pi}_{\hat H_\mu}=X^{2\pi}_{(\breve h+\breve f)^\R}$ there (sign $\sigma$ for the pair); in the rescaled slow variables $\tilde x=x/\lambda$ this is $X^c_{(\breve h+\breve f)^\R}$ on $W^e_r$.
\end{lemma}
\begin{proof}
\proofstep{Complex points}\label{stepN:p11:dm-pts} Let $(\theta,I,x)$ be a complex point with $\abs{\operatorname{Im}\theta}<2s=\frac\rho{2\pi}$, $\abs{I-I_c}<r\le\frac\rho2$, $\operatorname{dist}_\infty(x,S)<w$. Then $z=\bigl((2\pi\theta,q),(I,p)\bigr)\in\mathcal F(G',D,\rho)$: $\abs{\operatorname{Im}2\pi\theta_i}<\rho$; $\operatorname{Re}I$ is within $\frac\rho2$ of $I_c\in G$, so $\operatorname{Re}I\in G'$ (collar); $\abs{\operatorname{Im}I_i}<\frac\rho2$; there is $x_0\in S$ with $\abs{x-x_0}<w$, so $\operatorname{Re}x$ lies in the closed ball of radius $d/4$ about $x_0$, inside $D$ since $\bar B(x_0,3d/4)\subset D$; $\abs{\operatorname{Im}x_j}<w<\rho$. Hence $f_1$ is holomorphic there (P1), bounded by $\frac\mu{2\pi}$ (P2), $1$-periodic in $\theta$ (P4) and real at real points (P3). The same points show $V_r(I_c)\subset G'_\rho$ ($\operatorname{Re}I\in G'$ within $\rho$ of $I_c$, i.e.\ $I\in G'_\rho$), so (a) follows from the hypotheses and Lemma~\ref{bpN:p11:hessian}.
\proofstep{Norm}\label{stepN:p11:dm-norm} Lemma~\ref{bpN:p4:sup-weighted}(b) (extended, Lemma~\ref{bpN:p11:ext-part4}) with $\sigma=s$: $\abs{f_1}\le\frac\mu{2\pi}$ on $T_{2s}$-points gives $\wn{f_1}rs\le\coth(\pi s)^n\frac\mu{2\pi}=\varepsilon$. The domain is $\Dom{r}{2s}^e$ in the rescaled variables ($\operatorname{dist}(\tilde x,\tilde S)<r\iff\operatorname{dist}(x,S)<w$).
\proofstep{Field}\label{stepN:p11:dm-field} $\mathcal W(I_c,r)\subset\{\mathrm{cx}\,z\in\mathcal F\}$ (Step~\ref{stepN:p11:dm-pts}), which is open; on the real points of this open set $H_\mu=\operatorname{Re}(h+\mu P\circ\mathrm{cx})$ (frozen definition), so $\hat H_\mu$ is real analytic, equals the real part of the holomorphic function $h_1+f_1$ and $X^{2\pi}$ of it is as stated (Lemma~\ref{bpN:p11:real-system}(a) for the real parts). Rescaling: Step below Definition~\ref{bpN:p11:ext-class}.
\end{proof}

\begin{theorem}[Stopped block stability (Pos93 Theorem 1 with slow parameters)]\label{bpN:p11:block}
In Definition~\ref{bpN:p11:block-setting} assume in addition
\begin{enumerate}[label=(B\arabic*)]
\item $\kappa=2\pi K_as\ge6$ ($K_a=K/\sqrt n$, $K=(\varepsilon_0/\varepsilon)^{1/(2n)}$);
\item $K\ge256\,k\,\rho\,R_0/d^2$;
\end{enumerate}
and put $r_\Theta=r_0/(AK)^n$, $T_{\mathrm P}=\dfrac{s\,m_1\,r_\Theta}{100\,M_1\,\varepsilon}\,e^{\kappa/6}$. Let $(\breve h,\breve f)=(\sigma h_1,\sigma f_1)$ be the perturbation pair, $I_0\in G$, $F:\R^{2n+2k}\to\R^{2n+2k}$ any function with $F=X^{2\pi}_{(\breve h+\breve f)^\R}$ on $\mathcal W(I_0,r_0)$, $J\ni0$ an interval, $y=(\theta,I;x):\R\to\R^{2n+2k}$ differentiable on $J$ with $y'=F(y)$ there, $I(0)=I_0$, $x(0)\in S$, and $t_1\in J$ with $\abs{t_1}\le T_{\mathrm P}$ and $x(u)\in\operatorname{core}(D,d)$ for all $u$ between $0$ and $t_1$. Then $\abs{I(u)-I_0}_\infty\le\frac{R_0}{2K}=\frac{r_0}{2AK}$ for all $u$ between $0$ and $t_1$.
\source{Pos93 Theorem 1 (pp.~190--191, proof pp.~203--204) with Part~\ref{bpN:p9} (Theorem~\ref{bpN:p9:thm1}, Proposition~\ref{bpN:p9:block-hyp}) and the extension of this part; for the hypotheses of Part~\ref{bpN:p9} we take $P=\{I_0\}$, $f=0$, $\ell=2M_1r_0$ (convexity gives quasi-convexity for every $\ell$, and $r_0\le\ell/(2M_1)$ holds with equality).}
\end{theorem}
\begin{proof}
We write $h$ for $\breve h=\sigma h_1$ and $f$ for $\breve f=\sigma f_1$; Lemma~\ref{bpN:p11:domain} holds for the pair.
\proofstep{Block}\label{stepN:p11:bl-block} The setting of Definition~\ref{bpN:p9:setting} with $P=\{I_0\}$, $f=0$, $(r_0,s,m_1,M_1,\ell=2M_1r_0,h)$ satisfies (S1)--(S4) (Lemma~\ref{bpN:p11:hessian}: (S1); convexity: (S2); (S3) with equality; (S4) trivially). $K\ge1$ since $\varepsilon\le\varepsilon_0$. Lemma~\ref{bpN:p8:covering} gives $\Lambda\in\mathbb M_K$ with $I_0\in D_\Lambda$; as in Step~\ref{stepN:p9:t1-a-choice} (exclusion of the exceptional branch), $\norm{\omega(I_0)}_2\ge\abs{\omega(I_0)}_\infty\ge g_1\ge\frac18m_1r_0$ (because $r_0\le8g/m$), so $\Lambda\ne\Z^n$ (Step~\ref{stepN:p9:t1-a-choice}). Put $r=r_\Lambda$, $\alpha=\alpha_\Lambda$: $r_\Theta\le r\le r_0/(AK)=R_0/K$ (Lemma~\ref{bpN:p8:covering}(d)).
\proofstep{Block hypotheses}\label{stepN:p11:bl-hyp} Proposition~\ref{bpN:p9:block-hyp}(a)--(f) (their proofs use $f$ only through $\varepsilon_f\le\varepsilon$): domains, $\varepsilon\le m_1r^2/2^{10}=\frac pq\frac{m_1r^2}{2^7}$, $r=\alpha/(p\sqrt nM_1K_a)$, nonresonance modulo $\Lambda$ on $Z_{K_a}$ and $\norm{\Pi_\Lambda\omega(I_0)}_2<\delta_\Lambda\le\frac18m_1r$, the Lipschitz bound $L=\sqrt nM_1$ with $K_aLr=\alpha/p$, and for $\Lambda=\Theta$ the smallness $\varepsilon\le\alpha r/(2^7qK_a)$. These are \ref{bpN:p6:cond-nr}, \ref{bpN:p6:cond-L} and $K_aLr\le\alpha/p$ of Part~\ref{bpN:p6} (for $\Lambda=\Theta$ with $\Lambda=\{0\}$), which is what Propositions~\ref{bpN:p11:nonresonant}, \ref{bpN:p11:resonant} assume.
\proofstep{(SL)}\label{stepN:p11:bl-sl} $64\pi k s r\le w^2$ is implied by (B2): $64\pi ksr\le64\pi ks\frac{R_0}K=16k\rho\frac{R_0}K\le\frac{d^2}{16}=w^2$ since $s=\frac\rho{4\pi}$ and $K\ge256k\rho R_0/d^2$. (This is the only place where $d$, hence $b_n$, enters.)
\proofstep{Extended system}\label{stepN:p11:bl-ext} With $\lambda=w/r$, $\tilde y=(\theta,I,x/\lambda)$: $\tilde y'=X^c_{(h+f)^\R}(\tilde y)$ at points of $W^e_r(I_0)$ where $y\in\mathcal W(I_0,r)$ (Lemma~\ref{bpN:p11:domain}(c)); $F$ transported is an extension of $X^c_{(h+f)^\R}$ from $W^e_r$. The hypotheses $x(u)\in\operatorname{core}(D,d)\subset S$ give $\tilde x(u)\in\tilde S$ ($\operatorname{core}(D,d)\subset\operatorname{core}(D,3d/4)$), $x(0)\in S$.
\proofstep{Case $\Lambda=\Theta$}\label{stepN:p11:bl-theta} Proposition~\ref{bpN:p11:nonresonant} with $(r,s,K_a,\alpha,L,p)\to(r_\Theta,s,K_a,\alpha_\Theta,\sqrt nM_1,\frac98)$: its hypotheses are Step~\ref{stepN:p11:bl-hyp}, (B1), (SL). $T_{\rm nr}=\frac{sr_\Theta}{5\varepsilon}e^{\kappa/6}\ge T_{\mathrm P}$ ($\frac{m_1}{M_1}\le1$). Conclusion $\abs{I(u)-I_0}\le r_\Theta/2\le R_0/(2K)$.
\proofstep{Case $\Lambda\ne\Theta,\Z^n$}\label{stepN:p11:bl-res} Proposition~\ref{bpN:p11:resonant} with $(r,s,K_a,\alpha,\delta,p,m,M)\to(r_\Lambda,s,K_a,\alpha_\Lambda,\delta_\Lambda,\frac98,m_1,M_1)$ and $\delta=\frac18m_1r$ (any $\delta\ge\norm{\Pi\omega(I_0)}_2$ with $\delta\le\frac18m_1r$ will do): $T_\Lambda=\frac{sm_1r_\Lambda}{100M_1\varepsilon}e^{\kappa/6}\ge T_{\mathrm P}$ since $r_\Lambda\ge r_\Theta$. Conclusion $\norm{I(u)-I_0}_2\le r_\Lambda/2\le R_0/(2K)$.
\end{proof}

\subsection{Assembly: proof of the frozen theorem}\label{bpN:p11:s-final}

\begin{proof}[Proof of Theorem~\ref{bpN:p0:partial}]
\proofstep{Data and constants}\label{stepN:p11:final-data} Fix the data and hypotheses; take $\sigma$ from \texttt{HasDefiniteFastPart}; constants of Definition~\ref{bpN:p11:constants} (they depend only on the data). $G\subset G'$ ($I\in\bar B(I,\rho)\subset G'$) and $\mathcal F(G,D,\rho)$ is not needed.
\proofstep{The threshold}\label{stepN:p11:final-mu} With $\gamma_1,\gamma_2$ of Lemma~\ref{bpN:p11:exponents} and $c_K,\kappa_C,\varepsilon_0,R_0$ of Definition~\ref{bpN:p11:constants} let $\mu_0=\min\{\mu_a,\dots,\mu_f\}$,
\[\mu_a=\frac{2\pi\varepsilon_0}{\kappa_C},\ \mu_b=\rho^{1/b_n},\ \mu_c=\Bigl(\frac{\pi s\,c_K}{3\sqrt n}\Bigr)^{2n}\ (\text{so }\kappa\ge6),\ \mu_d=\Bigl(\frac{c_K}{256k\rho R_0}\Bigr)^{1/\gamma_1},\ \mu_e=\Bigl(\frac{c_K}{R_0}\Bigr)^{1/\gamma_2},\]
and $\mu_f>0$ such that $T_{\mathrm P}(\mu)\ge\mu^{-1}\exp(\mu^{-a_n})$ for $0<\mu<\mu_f$ (Step~\ref{stepN:p11:final-time}). Then for $0<\mu<\mu_0$: $\varepsilon\le\varepsilon_0$ ($\mu\le\mu_a$), $d=\mu^{b_n}<\rho$ ($\mu\le\mu_b$, $d$ increasing in... $\mu^{b_n}\le\rho$ iff $\mu\le\rho^{1/b_n}$), $\kappa=2\pi Ks/\sqrt n=2\pi c_Ks\mu^{-1/(2n)}/\sqrt n\ge6$ ($\mu\le\mu_c$), (B2): $K=c_K\mu^{-1/(2n)}\ge256k\rho R_0\mu^{-2b_n}$ iff $\mu^{\gamma_1}\le c_K/(256k\rho R_0)$ ($\mu\le\mu_d$), and $\frac{R_0}{2K}\le\frac12\mu^{b_n}$ iff $\mu^{\gamma_2}\le c_K/R_0$ ($\mu\le\mu_e$).
\proofstep{Time}\label{stepN:p11:final-time} $T_{\mathrm P}=C_T\,\mu^{-1/2}\exp(c_2\mu^{-1/(2n)})$ with $C_T=\frac{s\,m\,r_0}{100\,M_{\C}\,(Ac_K)^n}\cdot\frac{2\pi}{\kappa_C}>0$ ($r_\Theta=r_0(Ac_K)^{-n}\mu^{1/2}$, $\varepsilon=\kappa_C\mu/2\pi$, $\frac{m_1}{M_1}=\frac m{M_{\C}}$) and $c_2=\frac{\pi sc_K}{3\sqrt n}>0$. Since $a_n<\frac1{2n}$ (Lemma~\ref{bpN:p11:exponents}), $c_2\mu^{-1/(2n)}-\mu^{-a_n}\ge\frac{c_2}2\mu^{-1/(2n)}$ for $\mu$ small, and $\mu\,T_{\mathrm P}\exp(-\mu^{-a_n})\ge C_T\mu^{1/2}\exp(\frac{c_2}2\mu^{-1/(2n)})\to\infty$ (\ref{bpN:p11:ext-rpow}); so $\mu_f$ exists, depending only on the data.
\proofstep{Given $\mu,P,C,z,t$}\label{stepN:p11:final-given} Let $0<\mu<\mu_0$, $P$ with \texttt{IsPartialPerturbation}\,$G'\,D\,\rho\,P$, $C>0$, $z$ with (Z1), (Z2), and $t\in[0,\min(C,\mu^{-1}e^{\mu^{-a_n}})]$ with the slow confinement on $[0,t]$. If $t=0$ the claim is $0\le\frac12\mu^{b_n}$. Let $t>0$; $t\le\mu^{-1}e^{\mu^{-a_n}}\le T_{\mathrm P}$ (Step~\ref{stepN:p11:final-time}) and $[0,t]\subset[0,C]$.
\proofstep{The orbit}\label{stepN:p11:final-orbit} By (Z1) and the definition of \texttt{IsSolutionOn} (Definition~\ref{bpN:p0:sol}), $z$ is differentiable on $[0,C]$ with $z'=X_{H_\mu}(z)$ and $H_\mu$ differentiable at $z(u)$; $I_0:=\mathrm{fast}(z(0)_2)\in\operatorname{core}(G,\mu^{b_n})\subset G$ by (Z2). Put $y=\Theta\circ z$ (Lemma~\ref{bpN:p11:orbit}); $y'=X^{2\pi}_{\hat H_\mu}(y)$ on $[0,C]$, where $X^{2\pi}_{\hat H_\mu}$ is, for the total function $\hat H_\mu=\frac1{2\pi}H_\mu\circ\Theta^{-1}$, an everywhere defined function $F$ which equals $X^{2\pi}_{(h_1+f_1)^\R}$ on $\mathcal W(I_0,r_0)$ (Lemma~\ref{bpN:p11:domain}(c)). The slow variables of $y$ lie in $\operatorname{core}(D,d)$ on $[0,t]$ by the hypothesis, $I(0)=I_0$, $x(0)\in\operatorname{core}(D,d)\subset S$.
\proofstep{Sign}\label{stepN:p11:final-sign} If $\sigma=+1$ apply Theorem~\ref{bpN:p11:block} to $y$, $J=[0,C]$, $t_1=t$. If $\sigma=-1$ put $\bar y(u)=y(-u)$, $\bar J=[-C,0]$: $\bar y'(u)=-F(\bar y(u))$ and $-F=X^{2\pi}_{(-h_1-f_1)^\R}$ on $\mathcal W$, i.e.\ the field of the pair $(\sigma h_1,\sigma f_1)$; apply Theorem~\ref{bpN:p11:block} to $\bar y$, $\bar J$ and $t_1=-t$ ($\abs{t_1}\le T_{\mathrm P}$, slow confinement between $0$ and $-t$ is that of $y$ on $[0,t]$).
\proofstep{Drift}\label{stepN:p11:final-drift} In both cases $\abs{I(t)-I_0}_\infty\le\frac{R_0}{2K}\le\frac12\mu^{b_n}$ ($\mu\le\mu_e$). $I(t)=\mathrm{fast}(z(t)_2)$, $I_0=\mathrm{fast}(z(0)_2)$ by the definition of $\Theta$: this is the claim of Theorem~\ref{bpN:p0:partial}.
\proofstep{Quantifier order}\label{stepN:p11:final-order} $\mu_0$ depends only on $(n,k,\rho,g,m,M_{\C})$ (through $s,r_0,A,\varepsilon_0,R_0,\kappa_C,c_K,\gamma_{1,2}$); neither on $G,G',D,h$ beyond these data nor on $\mu,P,C,z,t$. This is stronger than required.
\end{proof}

\subsection{Deviations, errata, doubts, forward references, remaining work}\label{bpN:p11:s-dev}

\begin{remark}[Deviations from the source and from earlier parts]\label{bpN:p11:deviations}
(1) \emph{Route.} Not Nek77's \S\S5--10 (Remark~\ref{bpN:p11:route}); the proof is P\"oschel-type with slow parameters; it gives time $T_{\mathrm P}\ge C_T\mu^{-1/2}e^{c_2\mu^{-1/(2n)}}$ and drift $\le\frac{R_0}{2c_K}\mu^{1/(2n)}$, i.e.\ exponents $\frac1{2n}$ larger than $a_n,b_n$; the frozen exponents follow by monotonicity (Steps~\ref{stepN:p11:final-time}, \ref{stepN:p11:final-drift}).
(2) \emph{The extended structure.} The bracket constant $\frac2e$ of Parts~\ref{bpN:p4}--\ref{bpN:p6} is replaced by $1$ (Lemma~\ref{bpN:p11:bracket-g}); the numerical checks are Step~\ref{stepN:p11:nf-num}. The slow bracket is dominated by the fast one by (SL), a condition on the slow width.
(3) \emph{$\omega$-free resonant time.} Lemma~\ref{bpN:p11:energy-free} replaces the factor $1/\Omega$ of Pos93 pp.~198--200 by a bound independent of $\omega$.
(4) \emph{Convex case.} Lemma 1 of Pos93 (convex) is used instead of Lemma 2; no $\ell$-conditions.
(5) \emph{Initial erosion.} Only $\mathrm{fast}(z(0)_2)\in G$ is used; the erosion by $\mu^{b_n}$ is not needed (and no erosion of $G$ at later times: the solution is assumed in $G$ and the proof does not use it).
(6) \emph{Slow erosion.} The statement's margin $d=\mu^{b_n}$ is used with $S=\operatorname{core}(D,\frac34d)$, $w=\frac d4$; the normal-form orbit may differ from the true one by $\le w/2^{10}$ in the slow variables.
(7) \emph{$2\pi$ and max norms.} Converted once (Lemma~\ref{bpN:p11:orbit}); the Euclidean Hessian bound is $M_{\C}$ (Lemma~\ref{bpN:p11:hessian}).
(8) \emph{Imaginary width.} $r_0\le\rho/2$, $2s=\rho/(2\pi)$, $w<\rho$: all analytic requirements are inside the open tube of width $\rho$; no boundary continuity of $P$ is used.
\end{remark}

\begin{remark}[Suspected defects in the sources]\label{bpN:p11:errata}
(1) \emph{Nek77}, p.~1 (and p.~3): the proofs of the lemmas of \S10 are not in the file (``will be given elsewhere''); every use of Lemmas 10.3, 10.5, 10.7 is a \sourcegap{} in the printed source (Remark~\ref{bpN:p11:route}); the ledger should record the missing Part II as an acquisition gap, not a defect. (2) \emph{Nek77 Theorem 4.4, p.~30:} the erosion is by $d$ in the Euclidean metric of the source and the confinement hypothesis on $(p,q)$ is on the whole $[0,C]$; the frozen stopped quantifier is a deliberate strengthening whose proof (here) uses the confinement only on $[0,t]$. (3) \emph{Pos93 pp.~198--200:} the time bound $\propto1/\Omega$ (Part~\ref{bpN:p7}); not an error, but non-uniform for unbounded $G$ (Remark~\ref{bpN:p11:omega-unbounded}). (4) The remaining defects of Pos93 recorded in Parts~\ref{bpN:p6}, \ref{bpN:p7}, \ref{bpN:p9} (e.g.\ the radius hypothesis $Mr\le\ell/2$) do not occur here because the convex Lemma 1 is used.
\end{remark}

\begin{remark}[Doubts about the frozen statement]\label{bpN:p11:doubts}
No inconsistency of the frozen statement was found. Points to be aware of: (a) $\operatorname{core}(D,\mu^{b_n})$ uses closed max-norm balls in $\mathrm{Ph}_k$ (consistent with the polydisc used here); (b) \texttt{IsSolutionOn} only constrains $z$ on $[0,C]$ in $\{\mathrm{fast}\in G\}$; the slow variables may leave $D$ at times $>t$ with $P$ undefined there; this is harmless (Remark~\ref{bpN:p0:partial-totality}(c)); (c) the hypothesis $\bar B(I,\rho)\subset G'$ is used exactly twice: $V_r(I_c)\subset G'_\rho$ and $\operatorname{Re}I\in G'$ for the domain of $P$ (Lemma~\ref{bpN:p11:domain}); (d) the hypothesis $k\ge1$ is not used (the proof works for $k=0$, where $S=\{\mathrm{pt}\}$ and (SL) is empty); (e) $n\ge2$ is used in the covering (Lemma~\ref{bpN:p9:m-le-M}, $A^{2n}\ge A^2$) and in $a_n$.
\end{remark}

\begin{remark}[Forward references, status, remaining work]\label{bpN:p11:status}
\emph{Forward references:} none to unwritten parts; Part~\ref{bpN:p10}'s existence glue is not needed (the statement is stopped, orbits are supplied). \emph{Items to audit (new mathematics, not in a pinned source):} Lemma~\ref{bpN:p11:bracket} (slow bracket, elementary), Proposition~\ref{bpN:p11:lie} (Part~5 for the extended structure: the only new ingredients are the matrix $\Omega_c$, the constant $1$ for $\frac2e$, and the slow displacement), Theorem~\ref{bpN:p11:nf} (Step~\ref{stepN:p11:nf-num}: the two places where $\frac2e\le1$ was used), Lemma~\ref{bpN:p11:energy-free} (the $\omega$-free energy error; the constants $0.0507$, $0.0115$, $0.0702$, $0.1464$), Lemma~\ref{bpN:p11:transfer} (continuation with slow variables). \emph{Remaining work:} none for the blueprint; the proofs of Part~\ref{bpN:p11}b, Proposition~\ref{bpN:p11:lie}, are written as lists of changes to Parts~\ref{bpN:p5}--\ref{bpN:p6} (each change is a numbered step), and the Lean prover should first generalise Parts~\ref{bpN:p4}--\ref{bpN:p6} to a base set $\hat D\subset\C^{n+2k}$ with the matrix $\Omega_c$, so that Part~\ref{bpN:p11} becomes a corollary; an estimate for the Lean work: $2$k--$3$k lines for the generalisation of Parts~4--6 and $2$k lines for Lemmas~\ref{bpN:p11:energy-free}--\ref{bpN:p11:block} and the assembly.
\end{remark}
\endgroup
